\documentclass{article}
\title{PCA}
\author{QwanxingxuA}

\usepackage{ctex}
\usepackage{amsmath}
\usepackage{amssymb}

\begin{document}
\maketitle
    \section{PCA}
    PCA即主成分分析，它是使用正交变换将线性相关的变量转换为线性无关的变量。在数据处理中，通常使用相关性热力图来看变量间的线性相关性，对于
    强相关的变量我们会采取诸如合并、舍去其一等方法优化数据，但是显然此类方法是粗略的。

    何为主成分？字面意思就是主要的成分，换言之就是无线性相关的数据，因为它们是主要的。

    如何分析呢？理解主成分过后很简单了，正交变换过后的数据特征一定线性无关。但是这一映射过程
    会出现信息压缩或者说损失，我们不希望信息丢失太多。怎么分析这种信息问题，在SVD里面的分解借用奇异值，
    实际上奇异值是特征值的变形，特征值是很有用的，可以以特征值为桥梁搭建联系。这就是PCA。
    
    笔者在SVD那一章写道奇异值分解降维，实际上与PCA是一致的。后续笔者将详细介绍，并从SVD和PCA探索它们在降维上面到底干了什么。

    定义输入$\boldsymbol{x}\in\mathbb{R}^m$,输出$\boldsymbol{y}\in\mathbb{R}^m$。定义模型
    \[y_k=\boldsymbol{\alpha}_k^T\boldsymbol{x},\quad k=1,2,\cdots,m\]
    \[\boldsymbol{\alpha}_i\boldsymbol{\alpha}_j=0,\quad i\neq j,i=1,2,\cdots,m\quad j=1,2,\cdots,m\]
    \[\boldsymbol{\alpha}_k^T\boldsymbol{\alpha}_k=1\]
    
    \section{变量}
    PCA的操作往往都是针对坐标轴上生硬的数据执行的，这会导致使用者浅显的数据直觉。实际上，PCA针对的对象是随机变量，所以PCA仍然
    逃不开机器学习的概率范式。在前几章的笔记中，笔者特别强调概率范式下的模型定义：
    \[f:\boldsymbol{x}\rightarrow \boldsymbol{y}\]

    这是总体随机变量$\boldsymbol{x}$的模型。这意味着还有随机变量$\boldsymbol{x}_i$的模型
    \[f:\boldsymbol{x}_i\rightarrow y\]

    实际上这就是总体主成分分析和样本主成分分析模型，它们并不完全一样。

    笔者强调一点，坐标轴依旧是随机变量，不能理解成一个数，而是坐标轴向量$\boldsymbol{x}=\boldsymbol{e}_i$。

    \section{总体主成分分析}
    既然是考虑线性相关问题，自然可以纳入成熟的数学工具。

    协方差矩阵
    \[\Sigma=\operatorname{Cov}(\boldsymbol{x})=\mathbb{E}(\boldsymbol{x}-\mathbb{E}\boldsymbol{x})(\boldsymbol{x}-\mathbb{E}\boldsymbol{x})^T\]


    相关矩阵
    \[\boldsymbol{R}=\left[\frac{\operatorname{Cov}(x_i,x_j)}{\sqrt{\operatorname{Var}(x_i),\operatorname{Var}(x_j)}}\right]_{m\times m}\]

    采用怎样的策略呢。正交变换方式是很多的，PCA希望保留更多信息。所以按照信息保留程度选取模型。直观的来看，数据映射过后距离拉得越开，信息保留的往往越多，而拉得开可以使用
    方差：
    \[\operatorname{Var}(\boldsymbol{x})=\mathbb{E}(\boldsymbol{x}-\mathbb{E}\boldsymbol{x})^2=\Sigma\]
    
    于是算法迎刃而解了。定义$\mathcal{W}$为已经选取的模型。
    \begin{align*}
        \max_{\alpha_k} &\operatorname{Var}(y_k)=\operatorname{Var}(\alpha_k^T\boldsymbol{x})=\boldsymbol{\alpha}_k^T\boldsymbol{\Sigma}\boldsymbol{\alpha}_k\\
        s.t. &\boldsymbol{\alpha}_k^T\boldsymbol{\alpha}_k=1,\\
        &\boldsymbol{\alpha}_k\boldsymbol{\alpha}_i=0,\quad \boldsymbol{\alpha}_i\in \mathcal{W}
    \end{align*}

    解决这个优化问题比较简单。通过归纳法构造过拉格朗日函数求导即可，具体可以参考共轭梯度法。当然共轭梯度法比这个构造要复杂。但是解决了共轭构造，这类问题就不算难了。
    笔者不再赘述，给出结论：
    \[\boldsymbol{\Sigma}\boldsymbol{\alpha}_k-\lambda_k\boldsymbol{\alpha}_k\]
    其中，$\lambda_k$为第$k$步的拉格朗日乘子。

    可以发现，$\lambda_k$就是$\boldsymbol{\Sigma}$的特征值，对应的单位特征向量为$\boldsymbol{\alpha}_k$。

    那么令$\boldsymbol{A}=[\boldsymbol{\alpha}_1\quad\boldsymbol{\alpha}_2\quad\dots\quad \boldsymbol{\alpha}_m]$，有
    \[\boldsymbol{\Sigma}\boldsymbol{A}=\boldsymbol{A}\boldsymbol{\varLambda }\]
    \[\boldsymbol{\varLambda }=\operatorname{diag}(\lambda_1,\lambda_2,\cdots,\lambda_k)\]

    由于
    \[\boldsymbol{A}\boldsymbol{A}^T=\boldsymbol{A}^T\boldsymbol{A}=\boldsymbol{I}\]

    所以
    \[\boldsymbol{\Sigma}=\boldsymbol{A}\boldsymbol{\varLambda }\boldsymbol{A}^T\]

    所以
    \[\operatorname{Cov}(\boldsymbol{y})=[\boldsymbol{\alpha}_k^T\boldsymbol{\Sigma}\boldsymbol{\alpha}_k]_{m\times m}=\boldsymbol{\varLambda }\]

    推导了这么多，总结一下到底是什么结论。

    第一，正交变换向量就是协方差矩阵的单位特征向量。
    \[\boldsymbol{y}_k=\boldsymbol{\alpha}_k\boldsymbol{x}\]

    第二，变化过后的方差矩阵（也可以说是协方差矩阵）就是由特征值组成的对角矩阵。
    \[\operatorname{Var}(y)_k=\lambda_k\]
    \[\lambda_1\ge\lambda_2\ge\cdots\ge\lambda_k\]

    所以通过特征向量实现主成分提取，通过特征值分析信息保留程度。这样就完成了机器学习任务！

    \section{PCA与SVD}
    二者的思想方法可以说一致，只是方向不同。但是它们的根源都是半正定实对称矩阵。
    SVD利用半正定实对称矩阵：
    \[\boldsymbol{A}^T\boldsymbol{A}\]
    PCA利用半正定实对称矩阵即协方差矩阵：
    \[\operatorname{Cov}(x)=\boldsymbol{\Sigma}\]

    再看正交矩阵构造。SVD构造奇异值矩阵和特征向量矩阵，然后构造$\boldsymbol{U}$,于是得到三层变换：
    正交旋转变换、坐标轴缩放、正交旋转变换。
    \[\boldsymbol{A}=\boldsymbol{U}\boldsymbol{\Sigma}\boldsymbol{V}^T\]

    PCA,通过构造正交矩阵得到最优化问题，解问题的过程发现本质是特征值和单位向量。
    \[\boldsymbol{\Sigma}\boldsymbol{\alpha}_k=\lambda_k\boldsymbol{\alpha}_k\]

    发现没有，PCA是一个正交旋转，而SVD是分解成三个变换过程！而这根源是两种方法选择的半正定实对称矩阵的不同，
    也可以说是分解的出发点不一样。  

    PCA和SVD最大的相同点，就是是对称矩阵的对角化。正因如此，它们本质是一样的，要理解这一点，先统一它们的随机变量都是样本随机变量。

    对于样本随机变量，样本方差是
    \[\operatorname{Var}(x)=\frac{1}{n-1}\sum_{i=1}^{n}\operatorname{Var}(x_i)\]

    将样本数据规范化:
    \[x=\frac{x-\overline{x} }{\sqrt{\operatorname{Var}(x)}}\]

    相关矩阵或者协方差矩阵（标准化过后是一个矩阵，为了统一说法，下说相关矩阵）。
    \[\boldsymbol{R}=\frac{1}{n-1}\boldsymbol{X}\boldsymbol{X}^t\]

    下使用SVD表达式：
    \[\boldsymbol{A}\boldsymbol{A}^T=\boldsymbol{U}\boldsymbol{\Sigma}\boldsymbol{V}^T\boldsymbol{V}\boldsymbol{\Sigma}^T\boldsymbol{U}^T=\boldsymbol{U}\boldsymbol{\Sigma}\boldsymbol{\Sigma}^T\boldsymbol{U}^T\]

    如果理解了SVD，就会发现
    $\boldsymbol{\Sigma}_{SVD}\boldsymbol{\Sigma}_{SVD}^T=\boldsymbol{\Sigma}_{PCA}$

    于是
    $\boldsymbol{A}\boldsymbol{A}^T$就是协方差矩阵，所以对应的特征向量矩阵$\boldsymbol{U}$就是正交变换的解。

    哈哈，笔者说了这么多。实际上是一直没说
    \[\boldsymbol{A}=\boldsymbol{X}\]
    反而证明它们是一致的。

    对于信息的考量了，笔者在SVD说了，SVD的截断分解具有PCA的效果。

    总结一下，标准化过后的样本PCA和SVD是一致的。但是，其余情况，并不完全一致。

    \section{标准化补充}
    实际上PCA标准化是很重要的，不进行这一步是会导致不同的结果的，这也是为什么标准化过后才与SVD一致。

    首先扯一个话题，SVD和PCA谁更大，当然这个不好比较。这是一个读者问我的，他的疑惑是PCA有不标准化的推导，而标准过后与SVD一致，是不是PCA更大呢？

    笔者觉得这个问题很有意思。量纲会对数据产生影响，不标准化的PCA结果可能不同。从几何角度看，因为量纲的影响会导致数据区域变形，比如极度扁平的椭圆。
    所以人为的标准化消除这种影响。

    笔者觉得硬要比较大小的话还是SVD大。PCA量纲问题还是因为它只是正交旋转，但是SVD具有坐标轴压缩变换。对于其它场景呢，笔者觉得依旧是SVD大，主要是因为SVD更加深入矩阵的几何本质。
\end{document}